\section{Conclusion}
This paper introduced Adaptive Range Sorting (ARS), a hardware-conscious spatial classification framework designed to address the growing impact of memory hierarchy behavior, branch prediction costs, and parallel scalability constraints in modern sorting systems.

Rather than approaching sorting as a comparison-driven decision process, ARS transforms the problem into one of spatial classification through monotonic key projection, histogram-based partitioning, and contention-free scatter execution.
This perspective enables the framework to leverage predictable control flow, contiguous memory movement, and hardware-efficient partitioning strategies that align naturally with contemporary multi-core architectures.

The paper presented the evolutionary development of the ARS lineage, tracing its progression from recursive spatial trees to the Aero architecture.
Through the introduction of a cache-resident quantile mapping table and division-free classification pipeline, Aero improves occupancy stability under skewed distributions while maintaining high-throughput parallel execution.
The work also explored several experimental architectural branches that investigated recursive refinement, hierarchical memory staging, adaptive execution policies, and streaming-oriented processing.

To support continuous data ingestion, the Exp D architecture extended the framework beyond conventional batch sorting through epoch-based micro-batching and concurrent spatial consolidation.
This demonstrated that the same spatial classification principles used for parallel sorting can also be applied to streaming workloads while maintaining bounded memory growth and stable ingestion throughput.

Theoretical analysis showed that recursive spatial partitioning bounds computational work to linear complexity under fixed-width key assumptions, while practical analysis highlighted the importance of memory hierarchy behavior, cache residency, and branch reduction in determining real-world performance.
Experimental evaluation across diverse workloads demonstrated competitive throughput, strong scalability characteristics, robustness under skewed distributions, and favorable microarchitectural behavior relative to modern comparison-based baselines.

The results presented throughout this work suggest that spatial classification represents a viable alternative execution model for high-throughput data processing.
Rather than viewing sorting exclusively through the lens of comparison-based decision trees, ARS demonstrates that hardware-conscious partitioning, adaptive classification, and spatial organization can provide an effective foundation for parallel and streaming systems operating on modern computing platforms.

More broadly, the exploratory architectures investigated throughout this study indicate that ARS should be viewed not as a single algorithm, but as a family of spatial processing architectures.
The Aero and Exp D implementations represent only two points within a larger design space encompassing adaptive execution, hierarchical memory optimization, distributed processing, and accelerator-oriented execution models.
As hardware continues to evolve toward increasing parallelism and memory complexity, spatial classification may provide a promising direction for the next generation of scalable data-processing systems.
